\section{Hierarchical Sea-Surface Modeling}

\subsection{Overall Framework}

The objective of the proposed framework is to generate physically interpretable sea-surface realizations under controllable environmental conditions while retaining a computational form suitable for radar echo simulation. The environmental parameter vector is denoted by
\begin{equation}
\boldsymbol{p}_{e}
=
\left\{
U_{10},\,\phi_{w},\,\Omega,\,H_{s},\,L_{x},\,L_{y},\,\Delta x,\,\Delta y
\right\},
\label{eq:environment_parameter_vector}
\end{equation}
where $U_{10}$ is the wind speed at a height of 10~m, $\phi_{w}$ is the wind direction, $\Omega$ denotes the inverse wave age, $H_{s}$ is the prescribed or normalized significant wave height, $L_x$ and $L_y$ define the horizontal domain, and $\Delta x$ and $\Delta y$ are the sampling intervals. A random variable set $\boldsymbol{\xi}$ controls the spectral phases and therefore produces independent realizations under the same environmental condition. The radar parameter vector is written as
\begin{equation}
\boldsymbol{p}_{r}
=
\left\{
f_{c},\,\theta_{i},\,\phi_{r},\,\mathcal P,\,B,\,N_{p}
\right\},
\label{eq:radar_parameter_vector}
\end{equation}
where $f_c$, $\theta_i$, $\phi_r$, $\mathcal P$, $B$, and $N_p$ denote the carrier frequency, incidence angle, radar azimuth, polarization, signal bandwidth, and pulse number, respectively. Separating $\boldsymbol p_e$, $\boldsymbol p_r$, and $\boldsymbol\xi$ allows the sea state, radar observation, and random realization to be varied independently during dataset generation.

\begin{figure*}[!t]
\centering
\fbox{%
\parbox[c][3.8cm][c]{0.94\textwidth}{%
\centering
\textbf{Placeholder for the overall framework figure}\\[4pt]
Left to right: environmental parameters and random phases $\rightarrow$
linear Elfouhaily surface $\rightarrow$ nonlinear Lie background $\rightarrow$
equal-energy directional short-wave reconstruction $\rightarrow$
steep-crest detection and localized curling transformation $\rightarrow$
triangular facets and radar echo generation.\\[4pt]
The final version should show the three surface states using the same spatial domain and mark the locally triggered curling region.
}}
\caption{Overall framework of the proposed hierarchical sea-surface and radar-echo simulation method.}
\label{fig:overall_framework}
\end{figure*}

The surface construction contains three successive levels. First, a linear random surface is synthesized from the Elfouhaily directional spectrum \cite{elfouhaily1997unified}. Let $S_E(\boldsymbol{k};\boldsymbol p_e)$ be the two-dimensional spectrum and $\varphi_{\boldsymbol k}(\boldsymbol\xi)$ be the random phase associated with wave vector $\boldsymbol k$. The linear elevation is expressed compactly as
\begin{equation}
\begin{aligned}
\eta_{L}(x,y)
={}&\mathcal F^{-1}\!\left\{
\sqrt{2S_E(\boldsymbol{k};\boldsymbol p_e)
\Delta k_x\Delta k_y}\right.\\[-2pt]
&\left.{}\times
\exp\!\left[j\varphi_{\boldsymbol k}(\boldsymbol\xi)\right]
\right\},
\end{aligned}
\label{eq:linear_surface_framework}
\end{equation}
where conjugate symmetry is imposed to obtain a real-valued surface. A wind-dependent nonlinear Lie operator $\mathcal L$ is then applied to introduce crest sharpening and crest--trough asymmetry while constraining the large-area slope statistics:
\begin{equation}
\eta_{\mathrm{NL}}
=
\mathcal L
\left[
\eta_L;\,U_{10},\phi_w
\right].
\label{eq:nonlinear_background_framework}
\end{equation}
This level provides the nonlinear background surface but retains a single-valued elevation representation.

Second, a prescribed short-wave band $\mathcal K_s$ is reorganized to form wind-wave ridges with a dominant propagation direction. Instead of directly adding an independent component, the original band is removed and replaced by an equal-energy directional realization:
\begin{equation}
\eta_{D}
=
\eta_{\mathrm{NL}}
-\mathcal B_s\!\left[\eta_{\mathrm{NL}}\right]
+\eta_{s}^{\mathrm{dir}},
\label{eq:directional_band_replacement_framework}
\end{equation}
subject to
\begin{equation}
\int_{\mathcal K_s}
S_{s}^{\mathrm{dir}}(\boldsymbol k)\,\mathrm d\boldsymbol k
=
\int_{\mathcal K_s}
S_{\mathrm{NL}}(\boldsymbol k)\,\mathrm d\boldsymbol k.
\label{eq:directional_band_energy_constraint}
\end{equation}

Here, $\mathcal B_s[\cdot]$ is a band-pass operator, $\eta_s^{\mathrm{dir}}$ is the reconstructed directional component, and $S_s^{\mathrm{dir}}$ and $S_{\mathrm{NL}}$ are the corresponding spectral densities. Equation~\eqref{eq:directional_band_energy_constraint} prevents repeated energy injection into the selected band while allowing its directional distribution to be controlled.

Third, a local crest score $Q(x,y)$ is computed from the normalized elevation, the slope along the dominant propagation direction, and the negative directional curvature. The curling center is selected as
\begin{equation}
\begin{aligned}
(x_c,y_c)
&=\mathop{\arg\max}_{(x,y)\in\mathcal D_c}Q(x,y),\\
Q(x,y)&=w_hQ_h+w_sQ_s+w_cQ_c,
\end{aligned}
\label{eq:crest_selection_framework}
\end{equation}
where $\mathcal D_c$ is the admissible crest region and $w_h$, $w_s$, and $w_c$ are weighting coefficients. A local coordinate system $(u,v,z)$ is established at $(x_c,y_c)$, with $u$ along the dominant propagation direction and $v$ along the crest. The final curled surface is represented parametrically as
\begin{equation}
\boldsymbol r_C(u,v)
=
\boldsymbol r_D(u,v)
+W(u,v)\,
\boldsymbol d_C
\left(u,v,z;\boldsymbol p_c\right),
\label{eq:localized_curl_framework}
\end{equation}
where $\boldsymbol r_D=[x_D,y_D,\eta_D]^{\mathrm T}$ is the directional-wave surface, $W(u,v)$ is a compact support function, $\boldsymbol d_C$ is the curling displacement, and $\boldsymbol p_c$ contains the local length, width, height, forward-extension, and curling parameters. Unlike $\eta_L$, $\eta_{\mathrm{NL}}$, and $\eta_D$, the final geometry is not restricted to the form $z=\eta(x,y)$; the parametric coordinates in \eqref{eq:localized_curl_framework} permit locally overturned facets.

The parametric surface is triangulated to form the mesh $\mathcal M_C$, which provides facet locations, areas, and normal vectors for the radar model. The complete generation chain can therefore be summarized as
\begin{equation}
\begin{aligned}
\left(\boldsymbol p_e,\boldsymbol\xi\right)
&\xrightarrow{\;\mathcal G_E\;}
\eta_L
\xrightarrow{\;\mathcal L\;}
\eta_{\mathrm{NL}}\\[-2pt]
&\xrightarrow{\;\mathcal D_s\;}
\eta_D
\xrightarrow{\;\mathcal C\;}
\mathcal M_C
\xrightarrow{\;\mathcal R(\boldsymbol p_r)\;}
\boldsymbol y,
\end{aligned}
\label{eq:complete_simulation_chain}
\end{equation}
where $\mathcal G_E$, $\mathcal D_s$, $\mathcal C$, and $\mathcal R$ denote spectral synthesis, directional short-wave reconstruction, localized curling, and radar echo generation, respectively, and $\boldsymbol y$ is the simulated radar data. This hierarchical formulation preserves the efficiency of spectral synthesis over most of the domain and confines the additional curling computation to a sparse, physically selected region.











\subsection{Nonlinear Background Sea Surface}
\label{subsec:nonlinear_background}

The background surface must preserve the prescribed directional spectrum while introducing the crest--trough asymmetry that is absent from a linear Gaussian realization.  We therefore first synthesize an Elfouhaily surface and then apply a dimensionally consistent second-order Riesz/Lie correction.  All operations are performed in the wind-aligned coordinate system, which makes the directional dependence explicit and avoids an artificial rotation of the resulting sea state.

Let $k=(k_x^2+k_y^2)^{1/2}$ and let $\varphi=\operatorname{atan2}(k_y,k_x)$. The directional elevation spectrum is written as~\cite{elfouhaily1997unified}
\begin{equation}
 \Psi_{\eta}(k,\varphi)
 =\frac{S(k)}{k}D(k,\varphi-\varphi_w),
 \qquad
 S(k)=\frac{B_l(k)+B_h(k)}{k^3},
 \label{eq:elfouhaily_2d}
\end{equation}
where $B_l$ and $B_h$ denote the gravity- and capillary-wave curvature spectra, respectively, and $\varphi_w$ is the wind direction.  The directional spreading function used here is
\begin{equation}
 D(k,\Delta\varphi)
 =\frac{1}{2\pi}\left[1+\Delta(k)\cos(2\Delta\varphi)\right],
 \label{eq:directional_spreading}
\end{equation}
where $\Delta(k)$ is the Elfouhaily directionality parameter.  Thus, wind speed $U_{10}$ and inverse wave age $\Omega_c$ determine both the radial spectrum and its angular distribution.

On a discrete wavenumber grid, the linear Fourier coefficients are sampled as
\begin{equation}
 \widehat{\eta}_L(\boldsymbol{k})
 =\left[\Psi_{\eta}(\boldsymbol{k})
 \Delta k_x\Delta k_y\right]^{1/2}\xi(\boldsymbol{k}),
 \label{eq:linear_spectral_sampling}
\end{equation}
where $\xi$ is a zero-mean unit-variance complex Gaussian variable.  Hermitian symmetry, $\widehat{\eta}_L(-\boldsymbol{k})= \widehat{\eta}_L^{*}(\boldsymbol{k})$, is imposed explicitly so that $\eta_L(\boldsymbol{x})$ is real.  A single set of random phases is retained throughout the subsequent nonlinear transformation; consequently, differences between the linear and nonlinear surfaces arise from the model rather than from different random realizations.

To construct the nonlinear correction, the wavenumbers parallel and normal to the wind are defined by
\begin{equation}
 \begin{bmatrix}k_a\\k_c\end{bmatrix}
 =
 \begin{bmatrix}
  \cos\varphi_w & \sin\varphi_w\\
  -\sin\varphi_w & \cos\varphi_w
 \end{bmatrix}
 \begin{bmatrix}k_x\\k_y\end{bmatrix}.
 \label{eq:wind_aligned_wavenumbers}
\end{equation}

The associated Riesz fields are then obtained from
\begin{equation}
 \widehat{h}_a=-\mathrm{i}\frac{k_a}{k}
 \widehat{\eta}_L\mathcal{B}_i,
 \qquad
 \widehat{h}_c=-\mathrm{i}\frac{k_c}{k}
 \widehat{\eta}_L\mathcal{B}_i,
 \label{eq:riesz_fields}
\end{equation}
where $\mathcal{B}_i$ is the input-band mask.  Because $k_a/k$ and $k_c/k$ are dimensionless, $h_a$ and $h_c$ retain the dimension of elevation.  This property is essential for a dimensionally valid second-order correction.

Following the second-order Lie/Creamer construction used for nonlinear sea surfaces~\cite{zhang2023secondorder}, three quadratic terms are formed:
\begin{align}
 \widehat{L}_a
 &=-\frac{k_a^2}{2k}\mathcal{F}\{h_a^2\},
 &
 \widehat{L}_{ac}
 &=-\frac{k_a k_c}{k}\mathcal{F}\{h_a h_c\},
 \nonumber\\
 \widehat{L}_c
 &=-\frac{k_c^2}{2k}\mathcal{F}\{h_c^2\}.
 \label{eq:lie_quadratic_terms}
\end{align}

The resulting second-order correction is
\begin{equation}
 \widehat{L}
 =\widehat{L}_a+\widehat{L}_{ac}+\widehat{L}_c.
 \label{eq:second_order_lie}
\end{equation}

Each term in \eqref{eq:second_order_lie} has the dimension of elevation in the Fourier representation.  No dimensional wind-speed multiplier is attached to the correction; wind-speed dependence is carried by the input spectrum and by the slope-moment closure introduced below.

The preliminary nonlinear spectrum is
\begin{equation}
 \widehat{\eta}_{N,0}
 =\widehat{\eta}_L
 +\beta\,\mathcal{P}_{H}
 \left\{\mathcal{B}_o\widehat{L}\right\},
 \label{eq:preliminary_nonlinear_spectrum}
\end{equation}
where $\beta$ is a dimensionless nonlinear-strength coefficient, $\mathcal{B}_o$ is the output-band mask, and $\mathcal{P}_{H}$ denotes Hermitian projection.  Unless stated otherwise, $\mathcal{B}_i$ and $\mathcal{B}_o$ are truncated at $4k_p$ and $8k_p$, respectively.  The two-to-one bandwidth ratio retains the quadratic interactions while suppressing unresolved high-wavenumber products.

The Lie operation redistributes slope energy and therefore does not, by itself, guarantee agreement with the prescribed spectral mean-square slopes (MSSs). We close these second-order moments against the same Elfouhaily spectrum used to generate the surface.  For any spectrum $\widehat{\eta}$, the along- and cross-wind MSSs are evaluated as
\begin{equation}
 m_a=\sum_{\boldsymbol{k}}k_a^2
 \left|\widehat{\eta}(\boldsymbol{k})\right|^2,
 \qquad
 m_c=\sum_{\boldsymbol{k}}k_c^2
 \left|\widehat{\eta}(\boldsymbol{k})\right|^2.
 \label{eq:directional_mss}
\end{equation}

A positive angular amplitude modifier is introduced as
\begin{equation}
 G(\varphi;U_{10})=
 \exp\!\left\{\frac{1}{2}\left[
 \lambda_0(U_{10})+
 \lambda_2(U_{10})\cos 2(\varphi-\varphi_w)
 \right]\right\}.
 \label{eq:mss_closure_modifier}
\end{equation}

Independent calibration realizations are first used to determine two positive MSS scale factors, $s_a(U_{10})$ and $s_c(U_{10})$, by comparing the median raw nonlinear MSSs with the corresponding integrals of $\Psi_{\eta}$.  The logarithms of these factors are represented by smooth low-order functions of $\log U_{10}$ and are then held fixed for all validation realizations.  For a new realization, $\lambda_0$ and $\lambda_2$ are solved so that the dressed moments equal $s_a m_{a,0}$ and $s_c m_{c,0}$, where $(m_{a,0},m_{c,0})$ are the raw nonlinear moments.  This angular dressing is positive and does not alter the Fourier phases.  Hence, Cox--Munk and other empirical slope models are used only as external validation references, rather than as fitting constraints.

The final background surface is expressed as
\begin{equation}
 \eta_N(\boldsymbol{x},t)
 =\operatorname{Re}\!\left\{
 \mathcal{F}^{-1}\!\left[
 G(\varphi;U_{10})\widehat{\eta}_{N,0}
 \mathrm{e}^{-\mathrm{i}\omega(k)t}
 \right]\right\},
 \label{eq:final_nonlinear_background}
\end{equation}
where $\omega(k)$ follows the gravity--capillary dispersion relation.  This construction preserves a controllable environmental spectrum and directional MSS while adding non-Gaussian crest geometry.  The MSS closure controls only second-order slope moments; local overturning geometry is introduced separately in the curling-wave stage described later.

\begin{figure}[!t]
 \centering
\fbox{\parbox[c][3.2cm][c]{0.92\columnwidth}{\centering Placeholder for a comparison of the linear Elfouhaily realization and the nonlinear Lie-transformed background surface, including representative profiles and along-/cross-wind slope distributions.}} \caption{Illustration of the nonlinear background-surface construction. The linear and nonlinear cases use identical random phases.}
 \label{fig:nonlinear_background_construction}
\end{figure}












\subsection{Directional Short-Scale Wind-Wave Construction}
\label{subsec:directional_short_waves}

The nonlinear background contains broadband short-wave energy, but its local wave crests do not necessarily exhibit the concentrated propagation direction required for the subsequent curling-wave construction. Directly superimposing an independently generated component would count the selected spectral energy twice. We therefore remove a resolved short-wave band from the nonlinear background and replace it with a random directional component of equal energy. This operation changes the organization of the band without introducing an artificial increase in its variance.

Let $\lambda_{s,\min}$ and $\lambda_{s,\max}$ denote the prescribed wavelength limits. The corresponding wavenumber interval is
\begin{equation}
 k_l=\frac{2\pi}{\lambda_{s,\max}},
 \qquad
 k_h=\frac{2\pi}{\lambda_{s,\min}}.
 \label{eq:short_wave_band_limits}
\end{equation}
A smooth radial band-pass window $W_s(k)$ is used to avoid ringing at the two band edges:
\begin{equation}
W_s(k)=
\begin{cases}
0, & k\leq k_l^- \text{ or } k\geq k_h^+,\\
\frac{1}{2}\!\left[1-\cos\!\left(\pi\frac{k-k_l^-}{k_l-k_l^-}\right)\right], & k_l^-<k<k_l,\\
1, & k_l\leq k\leq k_h,\\
\frac{1}{2}\!\left[1+\cos\!\left(\pi\frac{k-k_h}{k_h^+-k_h}\right)\right], & k_h<k<k_h^+,
\end{cases}
\label{eq:smooth_short_wave_window}
\end{equation}
where $k_l^-=(1-\epsilon_b)k_l$, $k_h^+=(1+\epsilon_b)k_h$, and $\epsilon_b$ controls the transition width. The removed component and the retained residual are
\begin{align}
 \eta_{s,r}(\boldsymbol{x})
 &=\mathcal{F}^{-1}\!\left\{W_s(k)\widehat{\eta}_N(\boldsymbol{k})\right\},
 \nonumber\\
 \eta_r(\boldsymbol{x})
 &=\eta_N(\boldsymbol{x})-\eta_{s,r}(\boldsymbol{x}).
 \label{eq:background_band_separation}
\end{align}

The replacement spectrum combines a log-wavenumber envelope and an angular concentration function. Let $\boldsymbol{e}_s=(\cos\psi_s,\sin\psi_s)$ be the prescribed propagation direction and define
\begin{align}
 k_0&=\sqrt{k_lk_h},
 &
 \sigma_k&=\frac{1}{4}\log\!\left(\frac{k_h}{k_l}\right),
 \nonumber\\
 \theta_a&=\cos^{-1}\!\left(
 \frac{|\boldsymbol{k}\cdot\boldsymbol{e}_s|}{k}\right).
 \label{eq:short_wave_shape_parameters}
\end{align}
The nonnegative spectral shape is then written as
\begin{align}
 \Phi_s(\boldsymbol{k})
 &=W_s(k)
 \exp\!\left[-\frac{\log^2(k/k_0)}{2\sigma_k^2}\right]
 \nonumber\\
 &\quad\times
 \exp\!\left[-\frac{\theta_a^2}{2\sigma_\theta^2}\right].
 \label{eq:directional_short_wave_shape}
\end{align}
where $\sigma_\theta$ determines the directional spreading. The use of $|\boldsymbol{k}\cdot\boldsymbol{e}_s|$ produces an axially symmetric spectrum, which is required for a real-valued instantaneous surface; the physical traveling direction is assigned separately through the temporal phase. This concentrated angular distribution is consistent with the observed directional sensitivity of short wind waves~\cite{sergievskaya2024shortwaves}.

A preliminary Hermitian spectrum is sampled as
\begin{equation}
 \widehat{\eta}_{s,0}(\boldsymbol{k})
 =\mathcal{P}_H\!\left\{\sqrt{\Phi_s(\boldsymbol{k})}\,\xi_s(\boldsymbol{k})\right\},
 \label{eq:short_wave_random_sampling}
\end{equation}
where $\xi_s$ is an independent unit-variance complex Gaussian field and $\mathcal{P}_H$ denotes Hermitian projection. Its amplitude is normalized using the energy of the removed band:
\begin{equation}
 \alpha_s=\left[
 \rho_s\frac{\displaystyle\sum_{\boldsymbol{k}}
 \left|W_s(k)\widehat{\eta}_N(\boldsymbol{k})\right|^2}
 {\displaystyle\sum_{\boldsymbol{k}}
 \left|\widehat{\eta}_{s,0}(\boldsymbol{k})\right|^2}
 \right]^{1/2},
 \label{eq:short_wave_energy_normalization}
\end{equation}
where $\rho_s$ is the replacement-energy ratio. Setting $\rho_s=1$ gives
\begin{equation}
 \sum_{\boldsymbol{k}}\left|\alpha_s\widehat{\eta}_{s,0}\right|^2
 =\sum_{\boldsymbol{k}}\left|W_s\widehat{\eta}_N\right|^2,
 \label{eq:equal_energy_replacement}
\end{equation}
so the reconstructed short-wave component and the removed background band have equal variance by Parseval's theorem.

To obtain unidirectional temporal propagation without violating Hermitian symmetry, the positive and negative wavenumber branches are advanced with opposite signed frequencies:
\begin{align}
 \widehat{\eta}_s(\boldsymbol{k},t)
 &=\alpha_s\widehat{\eta}_{s,0}(\boldsymbol{k})
 \nonumber\\
 &\quad\times\exp\!\left[-\mathrm{i}\,\omega(k)
 \operatorname{sgn}(\boldsymbol{k}\cdot\boldsymbol{e}_s)t\right],
 \nonumber\\
 \omega(k)&=\sqrt{gk}.
 \label{eq:directional_short_wave_evolution}
\end{align}
The sea surface after band replacement is therefore
\begin{equation}
 \eta_D(\boldsymbol{x},t)
 =\eta_r(\boldsymbol{x},t)
 +\operatorname{Re}\!\left\{
 \mathcal{F}^{-1}\!\left[
 \widehat{\eta}_s(\boldsymbol{k},t)
 \right]\right\}.
 \label{eq:directional_short_wave_surface}
\end{equation}
In the implementation, $\lambda_{s,\min}=1.5$~m, $\lambda_{s,\max}=5.0$~m, $\sigma_\theta=8^{\circ}$, $\epsilon_b=0.15$, and $\rho_s=1$ unless stated otherwise. The minimum wavelength is represented by at least six grid samples, preventing unresolved grid-scale oscillations from being interpreted as physical wind waves. The parameters may be varied to generate different dominant directions, ridge spacings, and directional spreads while retaining the energy supplied by the background sea state.

\begin{figure}[!t]
 \centering
 \fbox{\parbox[c][3.4cm][c]{0.92\columnwidth}{\centering Placeholder for the removed background band, the equal-energy directional replacement component, and the resulting nonlinear sea surface.}}
 \caption{Directional short-wave band replacement. The replacement redistributes the selected band in wavenumber and direction while preserving its energy.}
 \label{fig:directional_short_wave_replacement}
\end{figure}
